\documentclass[aspectratio=169, 10pt]{beamer}

% ── Theme ───────────────────────────────────────────────────────────────────
\usetheme{Madrid}
\usecolortheme{beaver}
\setbeamertemplate{navigation symbols}{}
\setbeamertemplate{footline}[frame number]

% ── Packages ────────────────────────────────────────────────────────────────
\usepackage{amsmath, amssymb, bm}
\usepackage{booktabs}
\usepackage{xcolor}
\usepackage{graphicx}
\usepackage{array}
\usepackage{colortbl}

% ── Colour palette ──────────────────────────────────────────────────────────
\definecolor{cambblue}{RGB}{163,193,173}
\definecolor{darkblue}{RGB}{0,50,100}
\definecolor{darkred}{RGB}{150,0,0}
\definecolor{lightgrey}{RGB}{240,240,240}
\definecolor{gold}{RGB}{180,140,0}
\definecolor{mygreen}{RGB}{0,120,60}

\setbeamercolor{title}{bg=darkblue,fg=white}
\setbeamercolor{frametitle}{bg=darkblue,fg=white}
\setbeamercolor{block title}{bg=darkblue,fg=white}
\setbeamercolor{block body}{bg=lightgrey,fg=black}
\setbeamercolor{alerted text}{fg=darkred}

% ── Macros ──────────────────────────────────────────────────────────────────
\newcommand{\LSR}{\mathrm{LSR}}
\newcommand{\RMSE}{\mathrm{RMSE}}

% ── Title page ──────────────────────────────────────────────────────────────
\title[Lucas Critique \& ML]{The Lucas Critique Examined in ML}
\subtitle{D200 Machine Learning in Economics}
\author{4815L}
\institute{University of Cambridge}
\date{March 2026}

% ════════════════════════════════════════════════════════════════════════════
\begin{document}

\begin{frame}
  \titlepage
\end{frame}

\begin{frame}{Outline}
  \tableofcontents
\end{frame}

% ════════════════════════════════════════════════════════════════════════════
\section{Introduction \& Research Question}

\begin{frame}{Motivation: The Lucas Critique (1976)}

  \begin{columns}[T]
    \begin{column}{0.55\textwidth}
      \textbf{The problem:}
      \begin{itemize}\setlength\itemsep{4pt}
        \item Reduced-form parameters are \alert{not policy-invariant}
        \item Agents update expectations under policy change $\Rightarrow$ model coefficients shift
        \item ``Current-generation models'' rendered useless given constantly changing policy regime
      \end{itemize}

      \vspace{6pt}
      \textbf{Classical responses:}
      \begin{itemize}\setlength\itemsep{2pt}
        \item \textbf{VARs} — model-agnostic, reduced theory dependence
        \item \textbf{TVP-VARs} — time-varying parameters
        \item \textbf{MSM / Hamilton filter} — latent regime structure
        \item \textbf{DSGE} — micro-founded but still vulnerable
      \end{itemize}
    \end{column}
    \begin{column}{0.42\textwidth}
      \begin{alertblock}{Research Question}
        Do regime-switching ML models better withstand the Lucas Critique than classical alternatives — and \textbf{why}?
      \end{alertblock}

      \vspace{6pt}
      \begin{block}{Why ML?}
        \begin{itemize}\setlength\itemsep{2pt}
          \item Model-agnostic, data-driven selection
          \item Flexible functional form $\to$ lower approximation error
          \item But: overfitting may worsen post-break degradation
        \end{itemize}
      \end{block}
    \end{column}
  \end{columns}

\end{frame}

% ════════════════════════════════════════════════════════════════════════════
\section{Literature Review}

\begin{frame}{Literature: What Is Already Known}

  \begin{columns}[T]
    \begin{column}{0.48\textwidth}
      \textbf{Classical \& statistical context:}
      \begin{itemize}\setlength\itemsep{3pt}
        \item \textbf{Box (1976)} — ``all models are wrong but some are useful''
        \item \textbf{Breiman (2001)} — two cultures: data vs.\ algorithmic modelling
        \item \textbf{TVP-VARs} (Primiceri, 2005) — time-varying parameters
        \item \textbf{DSGEs} — structural; still misspecify when regime shifts (Storm, 2021)
      \end{itemize}
    \end{column}
    \begin{column}{0.48\textwidth}
      \textbf{ML and the Lucas Critique:}
      \begin{itemize}\setlength\itemsep{3pt}
        \item \textbf{Athey (2018)} — ML selects functional form / latent state from data
        \item \textbf{Chen (2018)} — Progresa / Malawi: ML strength when structure changes
        \item \textbf{Salle (2015)} — ANN expectations in agent-based model; no structural form
        \item \textbf{Pereira (2025)} — predictive economics methodology
      \end{itemize}
    \end{column}
  \end{columns}

  \vspace{8pt}
  \begin{block}{Gap This Paper Fills}
    No systematic head-to-head comparison of classical \emph{and} ML regime-switching models on the same Lucas Critique experiment — with both simulated (known DGP) and real macroeconomic data.
  \end{block}

\end{frame}

% ════════════════════════════════════════════════════════════════════════════
\section{Theory}

\begin{frame}{Theoretical Framework: Forecast Risk Decomposition}

  \textbf{Regime-dependent DGP:} $\;y_{t+1} = f_{s_t}(x_t) + \varepsilon_{t+1}$, $\quad s_t \in \{1,\dots,S\}$ (latent policy regime)

  \vspace{4pt}
  \textbf{MSE-optimal forecast:} $\;m^*(x_t) = \sum_s \Pr(s_t{=}s \mid x_t)\,\mathbb{E}[y_{t+1} \mid x_t, s_t{=}s]$

  \vspace{4pt}
  \[
    \mathbb{E}\!\left[(y_{t+1}-\hat{m})^2\right]
    = \underbrace{\sigma^2}_{\text{irreducible}}
    + \underbrace{\mathbb{E}\!\left[(m^*-m(\cdot;\mathcal{F}))^2\right]}_{\text{approximation error}}
    + \underbrace{\mathbb{E}\!\left[(m(\cdot;\mathcal{F})-\hat{m})^2\right]}_{\text{estimation error}}
  \]

  \vspace{2pt}
  \begin{columns}[T]
    \begin{column}{0.55\textwidth}
      \textbf{Under structural instability:}
      \begin{itemize}\setlength\itemsep{4pt}
        \item \textbf{Parametric models} — restrict $\mathcal{F}$; low estimation error but \alert{large approximation error} when $f_s$ is non-linear or regime assignment shifts
        \item \textbf{Flexible ML} — low approximation error pre-break; high \alert{estimation error} post-break (overfitting)
        \item \textbf{Variance-reducing break:} $\sigma^2\!\downarrow$ $\Rightarrow$ LSR $<1$ — \emph{inverse} Lucas Critique
      \end{itemize}
    \end{column}
    \begin{column}{0.42\textwidth}
      \begin{block}{Key Prediction}
        Neither class is immune. Models flexible in \textbf{latent-state estimation} should outperform those flexible only in functional form.
      \end{block}
      \vspace{4pt}
      \begin{block}{Ensemble Insight}
        Expanding $\mathcal{F}$ via model averaging reduces approximation error $\Rightarrow$ theoretical basis for ensemble stability.
      \end{block}
    \end{column}
  \end{columns}

\end{frame}

% ════════════════════════════════════════════════════════════════════════════
\section{Data}

\begin{frame}{Data: Two Experimental Settings}

  \begin{columns}[T]
    \begin{column}{0.5\textwidth}
      \textbf{Experiment A — Simulated DGP}

      Two-state Markov-switching AR(1) — known ground truth:
      \[
        y_t = \mu_{s_t} + \phi_{s_t} y_{t-1} + \sigma_{s_t}\varepsilon_t, \quad \varepsilon_t \overset{\text{iid}}{\sim}\mathcal{N}(0,1)
      \]

      \vspace{-2pt}
      \begin{center}\scriptsize
      \begin{tabular}{lcccc}
        \toprule
        & \multicolumn{2}{c}{Mild} & \multicolumn{2}{c}{Severe} \\
        \cmidrule(lr){2-3}\cmidrule(lr){4-5}
        Param & Pre & Post & Pre & Post \\
        \midrule
        $\mu_0$ & 0.5 & 0.2 & 0.5 & $-$0.5 \\
        $\mu_1$ & $-$1.5 & $-$1.0 & $-$1.5 & $-$3.0 \\
        $\sigma_0$ & 1.0 & 1.3 & 1.0 & 2.0 \\
        $\sigma_1$ & 2.5 & 3.0 & 2.5 & 4.5 \\
        $P_{00}$ & 0.95 & 0.90 & 0.95 & 0.80 \\
        \bottomrule
      \end{tabular}
      \end{center}
      {\scriptsize Mild $\approx$ gradual tightening; Severe $\approx$ Volcker-scale shift}\\[4pt]
      {\scriptsize Regime durations: expansion $\approx$20 periods; recession $\approx$10 periods}
    \end{column}
    \begin{column}{0.47\textwidth}
      \textbf{Experiment B — US FRED data}

      \vspace{4pt}
      \begin{tabular}{lll}
        \toprule
        Target & Break & Date \\
        \midrule
        INDPRO & Great Moderation & Jan 1984 \\
        CPIAUCSL & Volcker disinflation & Oct 1979 \\
        \bottomrule
      \end{tabular}

      \vspace{6pt}
      \textbf{Covariates:} EFFR, UNRATE, USREC (NBER binary)

      \vspace{4pt}
      \textbf{Features:} log-differences $\times100$, lags, 5/20-period rolling mean \& std

      \vspace{4pt}
      \textbf{Protocol:}
      \begin{itemize}\setlength\itemsep{2pt}
        \item Train on pre-break only
        \item Evaluate on post-break (no re-training)
        \item Report pre-RMSE, post-RMSE, and LSR
      \end{itemize}
    \end{column}
  \end{columns}

\end{frame}

% ════════════════════════════════════════════════════════════════════════════
\section{Methodology}

\begin{frame}{Models \& Evaluation Metric}

  \begin{center}\small
  \begin{tabular}{llll}
    \toprule
    Model & Family & Regime ID & Forecast \\
    \midrule
    AR(2) / ARMA(2,1) & Linear baseline & None & AR / AR+MA \\
    \rowcolor{lightgrey}
    MSM & Classical reg.-switch & Hamilton filter (EM) & Regime-weighted AR \\
    \rowcolor{lightgrey}
    TAR & Classical reg.-switch & Threshold on $y_{t-1}$ & Per-regime OLS \\
    \rowcolor{cambblue!60}
    HMM + Ridge & Probabilistic ML & Baum-Welch EM & Per-regime Ridge \\
    \rowcolor{cambblue!60}
    Mixture of Experts & Probabilistic ML & Soft gating (EM) & Weighted Ridge \\
    \rowcolor{cambblue!60}
    ML Regime (XGB) & Nonlinear ML & KMeans $\to$ XGBClassifier & Per-regime XGBRegressor \\
    \rowcolor{cambblue!60}
    MSNN & Nonlinear ML & Markov chain + MLP (EM) & Per-regime MLP \\
    \midrule
    Model Average & Ensemble & — & Equal-weight mean \\
    \bottomrule
  \end{tabular}
  \end{center}

  \vspace{4pt}
  \begin{block}{Lucas Sensitivity Ratio (LSR) — key metric}
    \[
      \LSR = \frac{\RMSE_{\text{post}}}{\RMSE_{\text{pre}}}
      \qquad
      \begin{cases}
        \LSR > 1 & \text{Lucas-vulnerable: post-break forecast worse} \\
        \LSR < 1 & \text{Inverse Lucas: variance-reducing break}
      \end{cases}
    \]
  \end{block}

\end{frame}

\begin{frame}{Model Details}

  \begin{columns}[T]
    \begin{column}{0.48\textwidth}
      \textbf{HMM + Ridge}
      \begin{itemize}\setlength\itemsep{2pt}
        \item Baum-Welch EM on $\mathbf{O}_t = [y_t,\, y_{t-1},\, \bar{y}^{(5)},\, s^{(5)},\, \bar{y}^{(20)}]^\top$
        \item Viterbi decoding $\to$ per-regime Ridge regressor
        \item Forward filter at test time — no look-ahead; soft posteriors adapt to shifted emissions without re-training
      \end{itemize}

      \vspace{4pt}
      \textbf{Mixture of Experts}
      \begin{itemize}\setlength\itemsep{2pt}
        \item $K{=}2$ Ridge experts; logistic gating via EM
        \item Soft forecast: $\hat{y}_t = \sum_k \pi_k(x_t)\,\hat\beta_k^\top x_t$
        \item Partial regularisation; reduced sensitivity to regime misclassification
      \end{itemize}
    \end{column}
    \begin{column}{0.48\textwidth}
      \textbf{ML Regime (XGB)}
      \begin{itemize}\setlength\itemsep{2pt}
        \item KMeans $\to$ XGBClassifier $\to$ per-regime XGBRegressor
        \item Frozen KMeans boundaries brittle to distributional shift
        \item Piecewise-constant leaves memorise pre-break structure $\to$ high LSR
      \end{itemize}

      \vspace{4pt}
      \textbf{MSNN}
      \begin{itemize}\setlength\itemsep{2pt}
        \item Hamilton MSM with per-regime MLP experts\\($n_\text{feat}\!\to\!32\!\to\!16\!\to\!1$, ReLU, Adam)
        \item EM: forward-backward smoother $\to$ soft $\gamma_t(k)$; analytical $\hat{P}_{jk}$
        \item Markov chain + smoothed posteriors act as regularisers (analogous to HMM)
      \end{itemize}
    \end{column}
  \end{columns}

\end{frame}

% ════════════════════════════════════════════════════════════════════════════
\section{Results}

\begin{frame}{Results: RMSE Across All Experiments}

  \begin{center}\small
  \begin{tabular}{lcccccccccccc}
    \toprule
    & \multicolumn{3}{c}{Simulated Mild} & \multicolumn{3}{c}{Simulated Severe} & \multicolumn{3}{c}{IP 1984} & \multicolumn{3}{c}{CPI 1979} \\
    \cmidrule(lr){2-4}\cmidrule(lr){5-7}\cmidrule(lr){8-10}\cmidrule(lr){11-13}
    Model & Pre & Post & LSR & Pre & Post & LSR & Pre & Post & LSR & Pre & Post & LSR \\
    \midrule
    AR(2)      & 1.634 & 2.706 & 1.66 & 1.634 & 4.261 & 2.61 & 0.717 & 0.514 & 0.72 & 2.631 & 2.594 & 0.99 \\
    ARMA(2,1)  & 1.674 & 5.577 & 3.33 & 1.674 & 11.61 & 6.94 & 0.811 & 0.514 & 0.63 & 2.533 & 3.556 & 1.40 \\
    MSM        & 1.673 & 5.968 & 3.57 & 1.673 & 11.96 & 7.15 & 0.823 & 0.723 & 0.88 & 2.954 & 4.038 & 1.37 \\
    TAR        & 1.537 & 2.526 & 1.64 & 1.537 & 4.266 & 2.78 & 0.624 & 0.397 & 0.64 & 1.912 & 2.276 & 1.19 \\
    \rowcolor{cambblue!40}
    \textbf{HMM} & 1.696 & 2.559 & \textbf{1.51} & 1.696 & 3.892 & \textbf{2.29} & 0.648 & \textbf{0.378} & \textbf{0.58} & 2.040 & 2.311 & \textbf{1.13} \\
    MoE        & 1.339 & 2.628 & 1.96 & 1.339 & 4.742 & 3.54 & 0.619 & 0.407 & 0.66 & 1.791 & 2.992 & 1.67 \\
    \rowcolor{red!20}
    XGB        & 0.396 & 3.337 & 8.43 & 0.396 & 8.471 & 21.41 & 0.140 & 0.443 & 3.18 & 0.357 & 2.559 & 7.18 \\
    MSNN       & 0.613 & \textbf{2.493} & 4.07 & 0.613 & \textbf{3.391} & 5.53 & 0.326 & 0.493 & 1.51 & 1.852 & 3.068 & 1.66 \\
    \rowcolor{cambblue!20}
    Avg        & 1.222 & \underline{2.015} & 2.47 & 1.222 & 6.007 & 4.92 & 0.521 & 0.401 & 0.77 & 1.757 & \underline{2.268} & 1.29 \\
    \bottomrule
  \end{tabular}
  \end{center}

  \vspace{2pt}
  {\footnotesize LSR $=$ post/pre RMSE. \colorbox{cambblue!40}{Blue} = best LSR. \colorbox{red!20}{Red} = worst. \textbf{Bold post} = best absolute post-break RMSE.}

\end{frame}

\begin{frame}{Key Results \& Discussion}

  \begin{columns}[T]
    \begin{column}{0.48\textwidth}
      \textbf{1. Lucas critique valid in all but IP 1984:}
      \begin{itemize}\setlength\itemsep{3pt}
        \item IP 1984: variance-reducing break — $\sigma^2\!\downarrow$ as irreducible error
        \item 7/9 models \emph{improve} post-break; only XGB and MSNN degrade (pre-break overfitting)
        \item All other experiments: post-break errors higher across the board
      \end{itemize}

      \vspace{6pt}
      \textbf{3. Post-break RMSE $\neq$ LSR:}
      \begin{itemize}\setlength\itemsep{3pt}
        \item MSNN: best post-break RMSE in both simulated experiments — but LSR = 4.07 / 5.53
        \item Model average: lower post-break RMSE \emph{and} mild LSR in 2/4 cases
        \item Both metrics needed; neither alone is sufficient
      \end{itemize}
    \end{column}
    \begin{column}{0.48\textwidth}
      \textbf{2. No clear ML vs.\ classical divide — mechanism matters:}
      \begin{itemize}\setlength\itemsep{3pt}
        \item \textbf{HMM best LSR} (all 4): flexible \emph{latent-state estimation}, not functional form
        \item \textbf{XGB worst LSR} (all 4): frozen KMeans + piecewise-constant leaves; best pre-RMSE (0.396) inflates LSR to 8.43/21.41
        \item \textbf{MSM collapses} despite correct specification — Hamilton filter regime accuracy $75\%\to20\%$; worse than AR(2)
        \item AR(2): reasonable post-break RMSE \emph{and} low LSR scores
      \end{itemize}

      \vspace{2pt}
      \begin{block}{Key Insight}
        Latent-state flexibility (HMM) $\gg$ functional-form flexibility (XGB). \alert{Correct specification can be detrimental.}
      \end{block}
    \end{column}
  \end{columns}

\end{frame}

% ════════════════════════════════════════════════════════════════════════════
\section{Conclusions}

\begin{frame}{Conclusions}

  \begin{columns}[T]
    \begin{column}{0.54\textwidth}
      \begin{block}{Main Findings}
        \begin{enumerate}\setlength\itemsep{4pt}\small
          \item \textbf{ML does not resolve the Lucas Critique} — better OOS forecasts can still worsen under unforeseen regime change
          \item \textbf{Latent-state flexibility is key} — HMM best LSR in all 4 experiments; XGB worst despite lowest pre-break RMSE
          \item \textbf{Correct specification can be detrimental} — MSM degrades more than AR(2) in all cases
          \item \textbf{Break character matters} — variance-reducing breaks (IP 1984) produce an \emph{inverse} Lucas Critique for 7/9 models
          \item \textbf{LSR and post-break RMSE diverge} — both needed; MSNN and Model Average illustrate the trade-off
        \end{enumerate}
      \end{block}
    \end{column}
    \begin{column}{0.43\textwidth}
      \begin{alertblock}{Policy Implication}
        ML does \emph{not} make forecasting models policy-invariant. HMM-family models preferred where structural stability matters.
      \end{alertblock}

      \vspace{6pt}
      \begin{block}{Future Work}
        \begin{itemize}\setlength\itemsep{3pt}
          \item \textbf{Ensembling} — key insight: reduces approximation error; optimised stacking / BMA
          \item \textbf{Online updating} — where does model memory most improve resilience?
          \item Richer DGPs and more real-world series
          \item Attention-based sequence models
        \end{itemize}
      \end{block}
    \end{column}
  \end{columns}

\end{frame}

% ── References ───────────────────────────────────────────────────────────────
\begin{frame}[allowframebreaks]{References}
  \footnotesize
  \begin{thebibliography}{99}
    \bibitem{lucas1976} R.~E. Lucas, ``Econometric policy evaluation: A critique,''
      \emph{Carnegie-Rochester Conference Series on Public Policy}, \textbf{1}, 19 (1976).

    \bibitem{hamilton1989} J.~D. Hamilton, ``A new approach to the economic analysis of
      nonstationary time series and the business cycle,''
      \emph{Econometrica}, \textbf{57}(2), 357--384 (1989).

    \bibitem{bormetti2022} G.~Bormetti and F.~Corsi, ``A Lucas critique compliant SVAR
      model with observation-driven time-varying parameters,'' arXiv:2107.05263 (2022).

    \bibitem{primiceri2005} G.~E. Primiceri, ``Time varying structural vector
      autoregressions and monetary policy,'' \emph{Review of Economic Studies},
      \textbf{72}(3), 821--852 (2005).

    \bibitem{storm2021} S.~Storm, ``Cordon of conformity: Why DSGE models are not the
      future of macroeconomics,'' \emph{Int.\ J.\ Political Economy},
      \textbf{50}, 77 (2021).

    \bibitem{box1976} G.~E.~P. Box, ``Science and statistics,''
      \emph{Journal of the American Statistical Association}, \textbf{71}, 791 (1976).

    \bibitem{breiman2001} L.~Breiman, ``Statistical modeling: The two cultures,''
      \emph{Statistical Science}, \textbf{16}, 199 (2001).

    \bibitem{athey2018} S.~Athey, ``The impact of machine learning on economics,''
      in \emph{The Economics of AI}, University of Chicago Press, 507--547 (2018).

    \bibitem{chen2018} T.-S. Chen, \emph{Essays on Applied Economics with Machine Learning},
      PhD dissertation, Washington University in St Louis (2018).

    \bibitem{salle2015} I.~L. Salle, ``Modeling expectations in agent-based models,''
      \emph{Economic Modelling}, \textbf{46}, 130 (2015).

    \bibitem{pereira2025} M.~A. Pereira, ``Predictive economics: Rethinking economic
      methodology with machine learning,'' arXiv:2510.04726 (2025).
  \end{thebibliography}
\end{frame}

\end{document}
